% Options for packages loaded elsewhere
\PassOptionsToPackage{unicode}{hyperref}
\PassOptionsToPackage{hyphens}{url}
\documentclass[
]{article}
\usepackage{xcolor}
\usepackage{amsmath,amssymb}
\setcounter{secnumdepth}{5}
\usepackage{iftex}
\ifPDFTeX
  \usepackage[T1]{fontenc}
  \usepackage[utf8]{inputenc}
  \usepackage{textcomp} % provide euro and other symbols
\else % if luatex or xetex
  \usepackage{unicode-math} % this also loads fontspec
  \defaultfontfeatures{Scale=MatchLowercase}
  \defaultfontfeatures[\rmfamily]{Ligatures=TeX,Scale=1}
\fi
\usepackage{lmodern}
\ifPDFTeX\else
  % xetex/luatex font selection
\fi
% Use upquote if available, for straight quotes in verbatim environments
\IfFileExists{upquote.sty}{\usepackage{upquote}}{}
\IfFileExists{microtype.sty}{% use microtype if available
  \usepackage[]{microtype}
  \UseMicrotypeSet[protrusion]{basicmath} % disable protrusion for tt fonts
}{}
\makeatletter
\@ifundefined{KOMAClassName}{% if non-KOMA class
  \IfFileExists{parskip.sty}{%
    \usepackage{parskip}
  }{% else
    \setlength{\parindent}{0pt}
    \setlength{\parskip}{6pt plus 2pt minus 1pt}}
}{% if KOMA class
  \KOMAoptions{parskip=half}}
\makeatother
\setlength{\emergencystretch}{3em} % prevent overfull lines
\providecommand{\tightlist}{%
  \setlength{\itemsep}{0pt}\setlength{\parskip}{0pt}}
\usepackage{enumitem}
\renewcommand{\thesection}{\S\arabic{section}}
\renewcommand{\thesubsection}{\arabic{section}.\arabic{subsection}}
\renewcommand{\thesubsubsection}{\arabic{section}.\arabic{subsection}.\arabic{subsubsection}}
\makeatletter
\renewcommand{\@seccntformat}[1]{\ifcsname the#1\endcsname\csname the#1\endcsname\hspace{0.4em}\fi}
\renewcommand{\numberline}[1]{#1\hspace{0.6em}}
\makeatother
% BTFV PDF header styling
\usepackage{graphicx}
\usepackage{xcolor}
\usepackage{fancyhdr}
\usepackage{geometry}
\usepackage{titlesec}
\usepackage{lastpage}
\usepackage{titling}
\usepackage{hyperref}
\usepackage{fontspec}

\definecolor{BTFVBlue}{HTML}{03326C}
\definecolor{BTFVDark}{HTML}{0F172A}
\definecolor{BTFVGray}{HTML}{6B7280}

\IfFontExistsTF{TeX Gyre Heros}{\setmainfont{TeX Gyre Heros}}{\setmainfont{Latin Modern Sans}}
\IfFontExistsTF{TeX Gyre Heros Condensed}{\newfontfamily\btfvtitlefont{TeX Gyre Heros Condensed}[Weight=500]}{\newfontfamily\btfvtitlefont{TeX Gyre Heros}[Weight=500]}

\geometry{bottom=1.8in}

\pagestyle{fancy}
\fancyhf{}
\renewcommand{\headrulewidth}{0pt}
\setlength{\headheight}{40pt}
\setlength{\footskip}{34pt}
\fancyhead[L]{}
\fancyhead[R]{
\includegraphics[height=50pt]{templates/btfv/images/logo.png}}
\fancyfoot[L]{\hspace*{15pt}\raisebox{-10pt}{\textcolor{BTFVGray}{\small \thedate}}}
\fancyfoot[C]{\raisebox{-15pt}{\IfFileExists{templates/btfv/images/footer_decoration.png}{\includegraphics[width=\textwidth]{templates/btfv/images/footer_decoration.png}}{}}}
\fancyfoot[R]{\raisebox{-10pt}{\textcolor{BTFVGray}{\small Seite \thepage\ von \pageref{LastPage}}}}

\titleformat{\section}{\Large\bfseries\color{BTFVBlue}}{\thesection}{0.6em}{}
\titleformat{\subsection}{\large\bfseries\color{BTFVBlue}}{\thesubsection}{0.6em}{}
\titleformat{\subsubsection}{\normalsize\bfseries\color{BTFVBlue}}{\thesubsubsection}{0.6em}{}

\setlength{\parskip}{0.6em}
\setlength{\parindent}{0pt}

\hypersetup{
  colorlinks=true,
  linkcolor=BTFVBlue,
  urlcolor=BTFVBlue
}

\makeatletter
\newcommand{\btfvsubtitle}{}
\@ifundefined{subtitle}{
  \newcommand{\subtitle}[1]{\gdef\btfvsubtitle{#1}}
}{
  \let\btfvorigsubtitle\subtitle
  \renewcommand{\subtitle}[1]{\gdef\btfvsubtitle{#1}\btfvorigsubtitle{#1}}
}
\makeatother

\makeatletter
\renewcommand{\maketitle}{
  \fancyhead[R]{}
  \begin{titlepage}
    \thispagestyle{fancy}
    \vspace*{1.2cm}
    \begin{center}
      
\includegraphics[height=150pt]{templates/btfv/images/logo.png}\par
      \vspace{0.9cm}
      {\Huge\btfvtitlefont\bfseries\color{BTFVBlue}\thetitle\par}
      \ifx\btfvsubtitle\@empty\else
        \vspace{0.35cm}
        {\Large\color{BTFVGray}\btfvsubtitle\par}
      \fi
      \vspace{0.8cm}
      {\large\color{BTFVGray}Stand: \thedate\par}
    \end{center}
    \vfill
  \end{titlepage}
  \setcounter{page}{2}
  \fancyhead[R]{
\includegraphics[height=50pt]{templates/btfv/images/logo.png}}
}
\makeatother
\usepackage{bookmark}
\IfFileExists{xurl.sty}{\usepackage{xurl}}{} % add URL line breaks if available
\urlstyle{same}
\hypersetup{
  pdftitle={Testdokument BTFV},
  hidelinks,
  pdfcreator={LaTeX via pandoc}}

\title{Testdokument BTFV}
\usepackage{etoolbox}
\makeatletter
\providecommand{\subtitle}[1]{% add subtitle to \maketitle
  \apptocmd{\@title}{\par {\large #1 \par}}{}{}
}
\makeatother
\subtitle{Design-Referenz für das BTFV-Template}
\author{}
\date{01.01.2000}

\begin{document}
\maketitle

\clearpage\renewcommand{\contentsname}{Inhaltsverzeichnis}
\tableofcontents
\clearpage

\section{Allgemeine Bestimmungen}\label{allgemeine-bestimmungen}

Dieser Abschnitt prüft die grundlegende §-Nummerierung mit dem
BTFV-Template. Der Abschnitts-Header soll in BTFV-Blau (\#03326C)
erscheinen und das BTFV-Logo oben rechts sichtbar sein.

\subsection{Zweck des Dokuments}\label{zweck-des-dokuments}

\begin{enumerate}
\def\labelenumi{\arabic{enumi}.}
\tightlist
\item
  Dieses Dokument dient als Referenz-PDF für automatisierte Tests.
\item
  Es prüft Schriftart (TeX Gyre Heros), Farben, Logo, Datum und
  §-Nummerierung.
\item
  Das Datum in der Fußzeile soll automatisch aufgelöst werden.
\end{enumerate}

\subsection{Geltungsbereich}\label{geltungsbereich}

\begin{enumerate}
\def\labelenumi{\arabic{enumi}.}
\tightlist
\item
  Das Dokument gilt für alle Build-Umgebungen (lokal und CI).
\item
  Es soll auf macOS und Ubuntu identische Ausgaben erzeugen.
\end{enumerate}

\section{Textformatierung und Listen}\label{textformatierung-und-listen}

\subsection{Geordnete Listen}\label{geordnete-listen}

Normale nummerierte Liste:

\begin{enumerate}
\def\labelenumi{\arabic{enumi}.}
\tightlist
\item
  Erster Punkt
\item
  Zweiter Punkt
\item
  Dritter Punkt
\end{enumerate}

\subsection{Alphabetische Listen}\label{alphabetische-listen}

\begin{enumerate}[label=\alph*.]
  \item Erster alphabetischer Punkt
  \item Zweiter alphabetischer Punkt
  \item Dritter alphabetischer Punkt
\end{enumerate}

\subsection{Ungeordnete Liste}\label{ungeordnete-liste}

\begin{itemize}
\tightlist
\item
  Stichpunkt A
\item
  Stichpunkt B

  \begin{itemize}
  \tightlist
  \item
    Unterpunkt B1
  \item
    Unterpunkt B2
  \end{itemize}
\item
  Stichpunkt C
\end{itemize}

\section{Sonderzeichen und
Typografie}\label{sonderzeichen-und-typografie}

\subsection{Umlaute}\label{umlaute}

Dieser Abschnitt enthält deutsche Umlaute: ä, ö, ü, Ä, Ö, Ü, ß.

\subsection{Hervorhebungen}\label{hervorhebungen}

Text kann \textbf{fett}, \emph{kursiv} oder \texttt{inline-code}
formatiert werden.

\subsection{Fußzeile und Datum}\label{fuuxdfzeile-und-datum}

Das Datum in der PDF-Fußzeile soll dem Ausstellungsdatum entsprechen.
Die Seitenzahl erscheint rechts in der Fußzeile im Format „Seite X von
Y''.

\section{Abschlussbestimmungen}\label{abschlussbestimmungen}

\begin{enumerate}
\def\labelenumi{\arabic{enumi}.}
\tightlist
\item
  Dieses Testdokument wird bei jedem Build neu generiert.
\item
  Abweichungen vom Referenz-PDF weisen auf Regressionen hin.
\end{enumerate}

\end{document}
